% Nouveau EGA V (cote 25): archive cote 25, site page 165 (PDF page 166).
% typescript. Draft compared: none.
% AI-transcribed by Claude Opus 5.5 (claude-opus-5-5, adaptive thinking, effort high), Message Batches API, batch ega25b (msgbatch_01A7n5YuXubUPNK1zsepi8RG), 2026-11-13; not yet checked by eye.
\scanpage{25/165}
% NOTE: the typed text of the whole page is a very faint carbon and almost unreadable on this scan; most words could not be read with any certainty
\illegible{4}

\illegible{2}

\uline{\unsure{Formalisme différentiel et morphismes lisses}}

\unsure{Pages 14 et 15} des notes \unsure{précédentes}, \illegible{4} \unsure{Page}
\unsure{15 ligne 20} le mot ``\illegible{1}'' est \unsure{absurde}, c'est la \unsure{deuxième assertion} qui est la
plus précise et celle à laquelle on songe.

\unsure{La généralisation de prop.} à laquelle il est fait allusion pages \unsure{14}
et 15 s'obtient à l'aide du diagramme suivant, qui généralise celui de la
page \unsure{4} :
\diagram{
\begin{tikzcd}
 & \unsure{\Omega^1_{?}\otimes\mathcal{O}_X} \arrow[d] & & \\
\unsure{0} \arrow[r] & \unsure{\underline{N}_{?}} \arrow[r] & \unsure{\Omega^1_{X/S}\otimes\mathcal{O}_{?}} \arrow[r] & \unsure{\Omega^1_{X/Y}} \arrow[r] & \unsure{0} \\
 & \unsure{\Omega^1_{Y/S}\otimes\mathcal{O}_X} \arrow[u] \arrow[urr] & &
\end{tikzcd}
}
\illegible{3} diagramme \illegible{3} \illegible{4} diagramme \illegible{1}
\illegible{4} \unsure{le lemme de la page} \illegible{2}
\illegible{3} dans le diagramme \unsure{précédent}. \illegible{1}
\unsure{les conditions de ce lemme}, on est ramené à prouver que la section
\illegible{1} \unsure{(?,?)} \unsure{déduit de} \illegible{2} \unsure{i.e.} \illegible{1} \unsure{$\to \Omega^1_{X/S}\otimes\mathcal{O}_?$}
\illegible{2} \illegible{1} est celui qui \illegible{4}
\illegible{2} satisfait \illegible{3} \illegible{4} de la
page \unsure{15}. Or les \illegible{3} \illegible{1}
\[ \unsure{V^1 \;\;\; V'} \;\illegible{3} \]
\margin{\unsure{||}}\unsure{donnent une suite exacte} \illegible{2}
\[ \unsure{0 \to \mathrm{Hom}_{?}(?)} \to \illegible{1} \to \illegible{1} \to 0 \; ; \]
\unsure{D'autre part} on a des \unsure{isomorphismes canoniques}
\[ \illegible{2} \simeq \illegible{1} \quad , \quad \illegible{1} \simeq \illegible{1} \quad , \]
\illegible{3} \unsure{(cf.} \illegible{1}\unsure{)} \illegible{4}

\mspage{\hand{165}}
